\documentclass[10pt,a4paper]{amsart}
\usepackage{amsmath,amssymb,amsthm,graphicx}
\RequirePackage{amssymb,amsmath,amsthm,bm}
\RequirePackage{thmtools}
\RequirePackage[capitalize]{cleveref}
\providecommand{\tightlist}{%
  \setlength{\itemsep}{0pt}\setlength{\parskip}{0pt}}

%%%%% varphi instead of phi
\renewcommand{\phi}{\varphi}

\mathchardef\mhyph="2D 
\providecommand{\A}{\mathcal{A}}
\providecommand{\C}{\mathcal{C}}
\providecommand{\L}{\mathcal{L}}
\providecommand{\B}{\mathcal{B}}
\providecommand{\M}{\mathcal{M}}
\renewcommand{\S}{\mathcal{S}}
\providecommand{\ol}[1]{\overline{#1}}
\providecommand{\mbf}[1]{\mathbf{#1}}
\providecommand{\mf}[1]{\mathfrak{#1}}
\providecommand{\mrm}[1]{\mathrm{#1}}
\providecommand{\mc}[1]{\mathcal{#1}}
\providecommand{\ra}{\rightarrow}
\providecommand{\Ra}{\Rightarrow}
\providecommand{\La}{\Leftarrow}
\providecommand{\lr}{\leftrightarrow}
\providecommand{\hf}{\mathbb{HF}}
\providecommand{\spec}{DgSp}
\providecommand{\bSigma}{\bm\Sigma}
\providecommand{\bDelta}{\bm\Delta}
\providecommand{\bPi}{\bm\Pi}
\providecommand{\Pinf}[1]{\Pi^{\mathrm{in}}_{#1}}
\providecommand{\Sinf}[1]{\Sigma^{\mathrm{in}}_{#1}}
\providecommand{\dSinf}[1]{d\mhyph\Sigma^{\mathrm{in}}_{#1}}
\providecommand{\Dinf}[1]{\Delta^{\mathrm{in}}_{#1}}
\providecommand{\bp}{\bar p}
\providecommand{\bq}{\bar q}
\providecommand{\ba}{\bar a}
\providecommand{\bb}{\bar b}
\providecommand{\bc}{\bar c}
\providecommand{\bd}{\bar d}
\providecommand{\bx}{\bar x}
\providecommand{\by}{\bar y}
\providecommand{\bz}{\bar z}
\providecommand{\LO}{\mathrm{LO}}
\declaretheorem[name={Theorem},numberwithin=section]{theorem}
\declaretheorem[name={Proposition}, sibling=theorem]{proposition}
\declaretheorem[name={Lemma}, sibling=theorem]{lemma}
\declaretheorem[name={Fact}, sibling=theorem]{fact}
\declaretheorem[name={Claim}, numberwithin=theorem]{claim}
\declaretheorem[name={Definition},numberwithin=section,sibling=theorem, style=definition]{definition}
\declaretheorem[name={Remark}, numbered=no,style=definition]{remark}
\declaretheorem[name={Corollary}, sibling=theorem]{corollary}
\declaretheorem[name={Question},sibling=theorem]{question}
\newcommand{\T}{{\mathcal{T}}}
\newcommand{\SR}{SR}
\newcommand{\pSR}{pSR}
\newcommand{\Q}{\mathbb{Q}}
\newcommand{\Z}{\mathbb{Z}}
\newcommand{\K}{\mathcal{K}}
\newcommand{\N}{\mathcal{N}}
\renewcommand{\L}{\mathcal{L}}
\providecommand{\Ginf}[1]{\Gamma^{\mathrm{in}}_{#1}}
\DeclareMathOperator{\restrict}{\upharpoonright}
\title{An explicit minimal three-universal class of countable linear orders}
\author{David Gonzalez, Dino Rossegger}
\date{arxiv:2305.07126v1; CC BY 4.0}
\usepackage[UTF8,fontset=fandol]{ctex}
\begin{document}
\maketitle
\noindent\textit{Editorial note.} All linear orders are countable. The class is interpreted up to isomorphism, and J means a proper isomorphism-invariant subclass as specified by the universality definition.\par\medskip
\noindent\textit{Editorial note.} The construction chooses representatives by endpoint behavior and the finite-block decomposition; it does not assert an isomorphism classification of every input linear order.\par\medskip
\noindent\textit{Editorial note.} 2026-10-11材料收口：复用既有H2、来源许可和近邻判断；核对原命题与作者证明起止、所指上下文及装配范围齐备。未重证作者结论；沿用既有非重复依据。\par\medskip
\noindent\textit{Editorial scope.} One original exact result. Native construction and load-bearing local dependencies preserved; named classical Scott-analysis inputs explicit.\par\medskip
\noindent\textit{Allowed standard prerequisites (editorial).}\begin{itemize}
\item Countable linear orders and infinitary back-and-forth. The supplied interval criterion for asymmetric back-and-forth; the unique countable dense endpoint-free order Q; successor relation is Pi1 and fixed finite successor configurations admit their usual infinitary descriptions.
\item Scott-sentence preservation. A Pi\_alpha Scott sentence is preserved in the indicated alpha-back-and-forth comparison. The short target-specific minimality lemma is included.
\end{itemize}\medskip
In this article we give a comprehensive Scott analysis of the class of linear
orderings using three definitions of rank. Given a structure $\A$ the
\emph{(parameterless) Scott rank} $\SR(\A)$, is the least $\alpha$, such that
the automorphism orbits of all tuples in $\A$ are $\Sinf{\alpha}$ definable; the
\emph{parameterized Scott rank} $\pSR(\A)$, is the least ordinal $\alpha$ such
that there is a parameter $\bp\in A^{<\omega}$ such that $SR((\A,\bp))=\alpha$;
and the \emph{Scott sentence complexity} of $\A$, $SSC(\A)$ is the least of the
complexities $\Pinf{\alpha}$, $\Sinf{\alpha}$, $\dSinf{\alpha}$ such that $\A$
has a Scott sentence of one of these complexities. While the latter is not an

\subsection{Preliminaries}
We assume that the reader is familiar with the basic techniques of Scott
analysis and with linear orderings in general. For background we refer to
Montalb\'an's upcoming book~\cite{montalban2021} and the standard reference by
Ash and Knight~\cite{ash2000}. Let us highlight the following lemma
which we will use as basic tools in our proofs without reference.
\begin{lemma}[{\cite[Lemma 15.8]{ash2000}}]\label{lem:bfandpartitions}  Suppose $\A,\B\in\LO$. Then $\A\leq_1 \B$ if and only if $\A$ is infinite or at least as large as $\B$. For $\alpha>1$, $\A\leq_\alpha \B$ if and only if for every $1\leq\beta<\alpha$ and every partition of $\B$ into intervals $\B_0,\dots,\B_n$, with endpoints in $\B$, there is a partition of $\A$ into intervals $\A_1,\dots,\A_n$ with endpoints in $\A$, such that $\B_i\leq_\beta \A_i$.
\end{lemma}
Notice that this lemma particularly implies that $(\A,\ba)\leq_\alpha (\B,\bb)$
if and only if $(-\infty, a_0)\leq_\alpha (-\infty, b_0)$,
$(a_{|\ba-1|},\infty)\leq_\alpha (b_{|\bb-1|},\infty)$ and for all $0\leq
i<|\ba-1|$ $(a_i,a_{i+1})\leq_\alpha (b_i,b_{i+1})$~\cite[Lemma 15.7]{ash2000}.

\begin{definition}\label{def:alphauniversality}
  Let $K_0\subseteq K_1$ be isomorphism invariant classes of structures. The
  class $K_0$ is \emph{$\alpha$-universal} for $K_1$ if for every $\A\in K_1$
  there is $\B\in K_0$ such that $\A\leq_\alpha\B$. 
\end{definition}

\begin{theorem}\label{thm:3universal}
The following class is $3$-universal for the class of linear orderings
$K=\{\omega+k, k+\omega^*, k+\Z+k':k,k'\in\mathbb N\}\cup\{\omega+\omega^*\}\cup
\{(\sum_{i\in k} n_i + m_i\cdot\Q )+n_{k+1}:n_i,m_i\in\mathbb N,
n_i>m_{i-1},m_{i+1}\}\cup\{k:k\in\mathbb N\}$. There is no $3$-universal class
$J\subset K$.
\end{theorem}

\begin{lemma}\label{lem:universalrankinvariant}
  If $K_0$ is $\alpha$-universal for $K_1$, then $K_0$ contains all structures
  in $K_1$ with $\Pinf{\alpha}$ Scott sentences.
\end{lemma}
\begin{proof}
  Suppose otherwise and consider $\A\in K_1\setminus K_0$ with a $\Pinf{\alpha}$
  Scott sentence. By $\alpha$-universality, there is $\B\in K_0$ with
  $\A\leq_\alpha\B$. But then $\A\cong \B$ and thus $\A\in K_0$, a
  contradiction.
\end{proof}

\begin{proof}
Given a linear order $\L$ we can write $\L=W+\K+R$ with $W$ well ordered or
empty, $R$ reverse well ordered or empty, and $\K$ empty or without a greatest or least element.

The element in $K$ that $\L$ is $\leq_3$-below will depend on the sizes
of $W$, and $R$ and the order type of $\K$. We will deal with these cases
independently.

\emph{Case 1: $W$ is infinite and $R$ is finite.}
In this case, $W=\omega+\alpha$ for some ordinal $\alpha$, so $\L=\omega+\K'+k$
for some $k\in \mathbb{N}$ and $\K'$ is either empty or without greatest element.
We claim that $\L\leq_3\omega+k$.
It is sufficient to show that $\omega+\K'\leq_3\omega$.
The $\exists$-player always wins the corresponding back-and-forth game by
playing the following strategy. In the first round they play isomorphically on
the initial $\omega$. Because of that they only need to win the back-and-forth game for $\omega\leq_2 \omega+\K'$. Say the $\forall$-player picks elements such that
$\omega+\K'$ is partitioned into $n$ intervals $\A_i$, each of cardinality $m_i$ for
$i<n$. The $\exists$-player must pick a partition of $\omega$ such that for all intervals $\B_i$, each of
cardinality $k_i$, $\A_i\leq_1 \B_i$. This inequality holds if and only if
$m_i\geq k_i$. Given this inequality, the intervals formed by picking the first $n$ elements of  $\omega$ 
clearly form the desired partition.

\emph{Case 2: $W$ is finite and $R$ is infinite.} This case is analogous to the
first case.

\emph{Case 3: Both W and R are infinite.}
In this case, we can write $\L=\omega+\K'+\omega^*$.
We claim that $\L\leq_3 \omega+\omega^*$. The winning strategy for the
$\exists$-player is similar to the one for case 1. In the first play they play
isomorphically to reduce the game to showing that $\omega+\omega^*\leq_2
\omega+\K'+\omega^*$. They can then win this game by picking partitions with
smaller intervals, exactly as above.

\emph{Case 3: Both W and R are finite.}
There are two subcases with different behavior here.
In both we assume that $\K$ has no greatest and least element, as if it was
empty, then $\L$ would be finite and thus isomorphic to a structure in $K$.

\emph{Subcase 1: There are arbitrarily large successor chains in $\K$.} Here we
claim that $\K\leq_3\Z$, which gives that $\L\leq_3 k+\Z+k'$ for $k,k'\in\mathbb{N}$.
The $\exists$-player always wins the corresponding game by playing the following
strategy. Assume the $\forall$-player plays an ordered tuple with the distance
between least and greatest element $n$. Then the $\exists$-player plays a
successor chain of size $n$ in $\K$ such that $\K=\K_1+n+\K_2$ where both $\K_1$
and $\K_2$ are infinite. It is then sufficient to show that $\omega^*\leq_2
\K_1$ and $\omega\leq_2 \K_2$. In these games, similar to the cases above, the
$\exists$-player just needs to ensure that they play partitions of smaller size
then the $\forall$-player and this is possible by the same strategy.

\emph{Subcase 2: There is an $n\in\mathbb N$ such that no element in $\K$ has
$n$ successors.} In this case, we consider the block relation $\sim$ given by
$a\sim b$ if and only if $[a,b]$ and $[b,a]$ are finite. By assumption every
block is a linear order of size $n$ or less. As no block can be infinite, no two
blocks can be successors in the quotient $\K/{\sim}$. Furthermore, as there is
no least or greatest element, there are no least or greatest blocks in
$\K/{\sim}$. Therefore, $\K/{\sim}=\Q$. In total, this means that
$\K=\sum_{q\in\Q} i_n(q)$ where $i_n:\Q\to\{1,\cdots, n\}$. If there are only
finitely many blocks of size $n$ , we may write $\K$ as a sum
$\K_1+n+\cdots+n+\K_l$ where no block of size $n$ occurs in any of the $\K_i$.
Repeating this process as necessary, we can write $\K$ as $\K_1+n_1+\K_2+n_2+\dots
\K_l$ where each $\K_i$ is given by $\K_i=\sum_{q\in(b_i,t_i)} i_n(q)$ such
that $\{ q\in (b_i,t_i): i_n(q)=m_i\}$ is infinite for $m_i=\max{\text{range}(i_n\restrict (b_i,t_i))}$.
Furthermore, we have that $m_i>n_i$ and if $i>1$, $m_i>n_{i-1}$.

We now show that $\L=k+\K_1+n_1+\dots
\K_l+k'\leq_3 k+m_1\cdot \Q+n_1\dots m_l\cdot \Q+k'$.
It is enough to show that $\K_i=\sum_{q\in(b_i,t_i)}i_n(q) \leq_3 m_i\cdot\Q$.
The $\exists$-player always wins this game by playing the following strategy.
Every element the $\forall$-player plays in the first play is in a block of size
$m_i$. The $\exists$-player can respond by playing elements in the same position
in blocks of size $m_i$. The resulting intervals in $\K_i$ are of the form
$l_1+\sum_{q\in (b_i^j,t_i^j)}i_n(q)+l_2$ where $b_i^j> b_i$ and $t_i^j<t_i$.
The corresponding intervals in $m_i\cdot \Q$ are of the form $l_1+m_i\cdot
\Q+l_2$. Winning the game thus comes down to finding a winning strategy for the
game $m_i\cdot \Q\leq_2 \sum_{q\in (b_i^j,t_i^j)}i_n(q)$. Given a partition in
$\sum_{q\in (b_i^j,t_i^j)}i_n(q)$, the $\exists$-player needs to find a
partition in $m_i\cdot \Q$ such that all intervals are smaller than the intervals
in the partition picked by the $\forall$-player. They can do
that using the fact that the finite blocks are densely ordered and always of maximum size.

This shows that $K$ is $3$-universal. Standard arguments show that each
element of $K$ has a $\Pinf3$ Scott sentence. Therefore, by \cref{lem:universalrankinvariant},
there is no $J\subset K$ such that $J$ is $3$-universal.
\end{proof}

\begin{thebibliography}{9}
\bibitem{ash2000} C. Ash and J. Knight. \emph{Computable Structures and the Hyperarithmetical Hierarchy}. Studies in Logic and the Foundations of Mathematics144,2000.
\bibitem{AR20} R. Alvir and D. Rossegger. The complexity of Scott sentences of scattered linear orders. \emph{Journal of Symbolic Logic}85(3)(2020),1079--1101.
\bibitem{ash1986} C.J. Ash. Recursive labelling systems and stability of recursive structures in hyperarithmetical degrees. \emph{Transactions of the American Mathematical Society}298(2)(1986),497--514.
\bibitem{montalban2021} A. Montalban. \emph{Computable Structure Theory: Beyond the Arithmetic}. Draft,2021. \texttt{https://math.berkeley.edu/\string~antonio/CSTpart2\_DRAFT.pdf}.
\bibitem{GM23} David Gonzalez and Antonio Montalb'an. \emph{The $\omega$-Vaught's Conjecture}. Native bibliography entry supplied in flattened.bbl.
\bibitem{MonIntermediate} Antonio Montalb'an. Classes of structures with no intermediate isomorphism problems. \emph{Journal of Symbolic Logic}81 (2016),127--150.
\end{thebibliography}
\end{document}
